\begin{table}[htbp]
\centering
\caption{Border premium in non-cannabis sectors, New Mexico county-quarter panel}
\label{tab:noncannabis_placebo}
\footnotesize
\begin{tabular}{lccccc}
\toprule
Sector & Coef. & Implied \% & $p$ & $N$ & Counties \\
\midrule
Alcohol retail (NAICS 445310) & $-0.612$ & $-45.8\%$ & $0.614$ & 172 & 25 \\
Grocery stores (NAICS 4451) & $+0.407$ & $+50.3\%$ & $0.275$ & 529 & 32 \\
Total retail trade (NAICS 44-45) & $+0.099$ & $+10.4\%$ & $0.441$ & 544 & 32 \\
\bottomrule
\end{tabular}
\begin{minipage}{0.98\textwidth}
\vspace{0.3em}
\footnotesize
\textit{Notes:} Dependent variable is $\log(1+\text{gross receipts})$ per county-quarter, by NAICS sector.
Border$_c$ = 1 if county $c$'s centroid lies within 75 miles of the nonlegal-state line, with $\geq$1,000
nonlegal-state residents within that reach (this paper's standard definition). Log county population control
and quarter fixed effects included; county-clustered standard errors. Alcohol retail's smaller $N$ and county
count reflect New Mexico's tax-confidentiality suppression of small-cell receipts, which is more frequent for
this narrower 6-digit category than for the broader grocery and total-retail categories.
\end{minipage}
\end{table}
